\section{Stationary trees and a lattice comparison}\label{sec:stationary}
\begin{proposition}[Finite stationary non-arithmetic trees]\label{prop:stationary}
Fix $k\ge2$ and $\rho_1,\ldots,\rho_k>0$ with $\sum_i\rho_i=1$.
On $\{1,\ldots,k\}^{\N}$ with the corresponding product probability,
take the cylinders indexed by all finite words, including the empty word.
Let $T$ be their cylinder operator. Put
\[
 a_i=-\log\rho_i,\qquad H=\sum_{i=1}^k\rho_i a_i.
\]
Assume there is no $h>0$ such that every $a_i$ belongs to $h\mathbb Z$.
Then, as $t\downarrow0$,
\[
 D(t)\sim\ncount(T;t)\sim\frac1{Ht}.
\]
\end{proposition}
\begin{proof}
The masses are $m_w=\prod_{j=1}^{|w|}\rho_{w_j}$, and
$\theta=\max_i\rho_i<1$. Vertices with mass above $t$ have bounded depth,
so $D(t)<\infty$. Proposition~\ref{prop:compact} applies.

Let $F$ be the probability measure assigning mass $\rho_i$ to $a_i$,
combining equal atoms if necessary, and let
$U=\sum_{\ell\ge0}F^{*\ell}$ be its renewal measure, including the
unit atom at zero. For $g(x)=e^{-x}\one_{\{x>0\}}$ and
$z(x)=e^{-x}D(e^{-x})$, direct summation over words gives
\[
 z(x)=\int_{[0,\infty)}g(x-y)\,dU(y).
\]
Indeed, a word contributes probability $m_w=e^{-\sum_j a_{w_j}}$
to the renewal measure, cancelling the factor $e^{\sum_j a_{w_j}}$
in $g(x-\sum_j a_{w_j})$.

The increments have finite positive mean $H$ and are non-arithmetic
by hypothesis. The function $g$ is directly Riemann integrable.
The key renewal theorem therefore gives
\[
 z(x)\longrightarrow\frac1H\int_0^\infty e^{-u}\,du=\frac1H.
\]
For a precise source, see Gallager~\cite{Gallager}, Theorem~4.7.3,
printed page~202. That theorem uses the renewal measure without the
initial atom; the additional term here is $g(x)\to0$.
Thus $D(t)\sim1/(Ht)$.

All normalized subtrees are unitarily equivalent. Their shifts
$\widehat\xi_v$ are consequently one common function $h$.
It is bounded by one and supported in $[0,M_\theta]$, and has integral
one by Lemma~\ref{lem:rankone-shift}. Away from zero its discontinuities
are among finitely many eigenvalues on each compact subinterval; it is
therefore Riemann integrable on $[0,M_\theta]$. Apply
Theorem~\ref{thm:profile} with one class and zero error envelopes to get
$\ncount(T;t)\sim D(t)$.
\end{proof}

For example, the weights $(3/10,7/10)$ satisfy the non-arithmetic
condition. A rational relation between $\log(10/3)$ and $\log(10/7)$
would, after exponentiation and comparison of prime exponents, give a
nontrivial equality between powers involving distinct primes.

\begin{proposition}[Homogeneous $k$-ary trees]\label{prop:kary}
In the same construction with $\rho_i=1/k$ for all $i$, the spectrum is
\[
 \frac{k}{k-1}\quad\text{with multiplicity }1,\qquad
 \frac{k^{-\ell}}{k-1}\quad\text{with multiplicity }(k-1)k^\ell
 \quad(\ell\ge0).
\]
In particular,
\begin{equation}\label{eq:kary-count}
 \ncount(T;t)=\one_{\{t<k/(k-1)\}}+(k-1)D((k-1)t),
 \qquad
 D(t)=\sum_{\ell\ge0}k^\ell\one_{\{k^{-\ell}>t\}}.
\end{equation}
The function $D$ is not regularly varying at zero with index $-1$.
For $k=3$,
\[
 \liminf_{t\downarrow0}\frac{\ncount(T;t)}{D(t)}=\frac23,\qquad
 \limsup_{t\downarrow0}\frac{\ncount(T;t)}{D(t)}=2.
\]
\end{proposition}
\begin{proof}
Let $E_\ell$ be conditional expectation onto the first $\ell$ coordinates,
with $E_0$ the projection onto constants. The sum of cylinder rank-one
operators at depth $\ell$ is $k^{-\ell}E_\ell$. Thus
$T=\sum_{\ell\ge0}k^{-\ell}E_\ell$.
Constants have eigenvalue $\sum_{\ell\ge0}k^{-\ell}=k/(k-1)$.
The martingale difference space
$\operatorname{ran}(E_{\ell+1}-E_\ell)$ has dimension $(k-1)k^\ell$,
and $T$ acts on it by
$\sum_{j\ge\ell+1}k^{-j}=k^{-\ell}/(k-1)$.
These orthogonal spaces and the constants exhaust $L^2$.
This proves the spectrum and \eqref{eq:kary-count}.

For any fixed $c\in(1,k)$ and $t_m=k^{-m}$, the strict counting
convention gives $D(ct_m)=D(t_m)$. This is inconsistent with the ratio
$c^{-1}$ required for regular variation of index $-1$.

For $k=3$, in each mass interval the argument $2t$ crosses exactly one
mass level. The ratios $D(2t)/D(t)$ have limiting values $1/3$ and $1$,
by summing the geometric series in \eqref{eq:kary-count}.
The single constant eigenvalue is negligible as $t\downarrow0$,
giving the two stated limits.
\end{proof}

\begin{remark}
For $k=2$, \eqref{eq:kary-count} instead gives
$\ncount(T;t)=D(t)+1$ for $0<t<2$.
Thus a nonvanishing child ratio can coexist with exact asymptotic transfer.
The $k=3$ example shows that $\theta<1$ alone does not force transfer,
but it does not answer the question with regular variation retained:
its $D$ fails that hypothesis.
\end{remark}
